%%%%%%%%%%%%%%%%%
% This is an sample CV template created using altacv.cls
% (v1.7.4, 30 July 2025) written by LianTze Lim (liantze@gmail.com). Compiles with pdfLaTeX, XeLaTeX and LuaLaTeX.
%
%% It may be distributed and/or modified under the
%% conditions of the LaTeX Project Public License, either version 1.3
%% of this license or (at your option) any later version.
%% The latest version of this license is in
%%    http://www.latex-project.org/lppl.txt
%% and version 1.3 or later is part of all distributions of LaTeX
%% version 2003/12/01 or later.
%%%%%%%%%%%%%%%%

%% Use the "normalphoto" option if you want a normal photo instead of cropped to a circle
% \documentclass[10pt,a4paper,withhyper,normalphoto]{altacv}

\documentclass[10pt,a4paper,withhyper]{../altacv}
%% AltaCV uses the fontawesome5 and simpleicons packages.
%% See http://texdoc.net/pkg/fontawesome5 and http://texdoc.net/pkg/simpleicons for full list of symbols.

% Change the page layout if you need to
\geometry{left=1.25cm,right=1.25cm,top=1.5cm,bottom=1.5cm,columnsep=1.2cm}

% The paracol package lets you typeset columns of text in parallel
\usepackage{paracol}
\usepackage{xeCJK}

% Change the font if you want to, depending on whether
% you're using pdflatex or xelatex/lualatex
% WHEN COMPILING WITH XELATEX PLEASE USE
% xelatex -shell-escape -output-driver="xdvipdfmx -z 0" sample.tex

% using XeLaTex
% English font
\setCJKmainfont{Roboto Slab}
\setCJKsansfont{Lato}

% Chinese font
\setCJKmainfont{Microsoft YaHei}
\setCJKsansfont{Microsoft YaHei}

\renewcommand{\familydefault}{\sfdefault}


% Change the colours if you want to
\definecolor{SlateGrey}{HTML}{2E2E2E}
\definecolor{LightGrey}{HTML}{666666}
\definecolor{DarkPastelRed}{HTML}{450808}
\definecolor{PastelRed}{HTML}{8F0D0D}
\definecolor{GoldenEarth}{HTML}{E7D192}
\colorlet{name}{black}
\colorlet{tagline}{PastelRed}
\colorlet{heading}{DarkPastelRed}
\colorlet{headingrule}{GoldenEarth}
\colorlet{subheading}{PastelRed}
\colorlet{accent}{PastelRed}
\colorlet{emphasis}{SlateGrey}
\colorlet{body}{LightGrey}

% Change some fonts, if necessary
\renewcommand{\namefont}{\Huge\rmfamily\bfseries}
\renewcommand{\personalinfofont}{\footnotesize}
\renewcommand{\cvsectionfont}{\LARGE\rmfamily\bfseries}
\renewcommand{\cvsubsectionfont}{\large\bfseries}


% Change the bullets for itemize and rating marker
% for \cvskill if you want to
\renewcommand{\cvItemMarker}{{\small\textbullet}}
\renewcommand{\cvRatingMarker}{\faCircle}
% ...and the markers for the date/location for \cvevent
% \renewcommand{\cvDateMarker}{\faCalendar*[regular]}
% \renewcommand{\cvLocationMarker}{\faMapMarker*}


% If your CV/résumé is in a language other than English,
% then you probably want to change these so that when you
% copy-paste from the PDF or run pdftotext, the location
% and date marker icons for \cvevent will paste as correct
% translations. For example Spanish:
% \renewcommand{\locationname}{Ubicación}
% \renewcommand{\datename}{Fecha}


%% Use (and optionally edit if necessary) this .tex if you
%% want to use an author-year reference style like APA(6)
%% for your publication list
% \input{pubs-authoryear}

%% Use (and optionally edit if necessary) this .tex if you
%% want an originally numerical reference style like IEEE
%% for your publication list
\usepackage[backend=biber,style=ieee,sorting=ydnt,defernumbers=true]{biblatex}
%% For removing numbering entirely when using a numeric style
\setlength{\bibhang}{1.25em}
\DeclareFieldFormat{labelnumberwidth}{\makebox[\bibhang][l]{\itemmarker}}
\setlength{\biblabelsep}{0pt}
\defbibheading{pubtype}{\cvsubsection{#1}}
\renewcommand{\bibsetup}{\vspace*{-\baselineskip}}
\AtEveryBibitem{%
  \iffieldundef{doi}{}{\clearfield{url}}%
}



\begin{document}
\name{王冠麟}
\tagline{软件与 IT 工程师}
%% You can add multiple photos on the left or right
\photoR{2.8cm}{../../images/gl_photo}
% \photoL{2.5cm}{Yacht_High,Suitcase_High}

\personalinfo{%
  % Not all of these are required!
  \email{g.wang@mail.de}
  \phone{+49 151 56312173}
  \mailaddress{Alfred-Messel-Weg 38, 64287, Darmstadt}
  \printinfo{\faLinkedin}{Guanlin Wang}[https://www.linkedin.com/in/guanlin-wang-737354206]
  % \printinfo{\simpleicon{linkedin}\hspace{0.3em}Guanlin Wang}[https://www.linkedin.com/in/guanlin-wang-737354206]
}

\makecvheader
%% Depending on your tastes, you may want to make fonts of itemize environments slightly smaller
% \AtBeginEnvironment{itemize}{\small}

%% Set the left/right column width ratio to 6:4.
\columnratio{0.3}

% Start a 2-column paracol. Both the left and right columns will automatically
% break across pages if things get too long.
\begin{paracol}{2}

  \cvsection{个人能力}

  % Don't overuse these \cvtag boxes — they're just eye-candies and not essential. If something doesn't fit on a single line, it probably works better as part of an itemized list (probably inlined itemized list), or just as a comma-separated list of strengths.
  \cvsubsection{系统 \& 部署}
  \cvtag{Linux}

  \vspace{0.5em}

  \cvtag{Ubuntu / WSL}

  \vspace{0.5em}

  \cvtag{Docker}

  \vspace{0.5em}

  \cvtag{SSH}

  \vspace{0.5em}
  
  \cvtag{Shell / Bash}
  
  \vspace{0.5em}

  \cvtag{Batch}

  \vspace{0.5em}

  \divider\smallskip

  \cvsubsection{软件开发}
  
  \cvtag{Python}
  
  \vspace{0.5em}

  \cvtag{C++}
  
  \vspace{0.5em}

  \cvtag{C\#}
  
  \vspace{0.5em}

  \cvtag{.NET}
  
  \vspace{0.5em}

  \cvtag{REST API}
  
  \vspace{0.5em}

  \cvtag{FastAPI}
  
  \vspace{0.5em}

  \cvtag{ASP.NET Core}
  
  \vspace{0.5em}

  \divider\smallskip

  \cvsubsection{数据库 \& 集成}
  
  \cvtag{Microsoft SQL Server}
  
  \vspace{0.5em}

  \cvtag{SQL}
  
  \vspace{0.5em}

  \cvtag{System Integration}

  \vspace{0.5em}

  \divider\smallskip

  \cvsubsection{AI \& 图像处理}
  
  \cvtag{Image Processing}
  
  \vspace{0.5em}  
  
  \cvtag{OCR}
  
  \vspace{0.5em}  
  
  \cvtag{Computer Vision}
  
  \vspace{0.5em}  
  
  \cvtag{OpenCV}
  
  \vspace{0.5em}  
  
  \cvtag{PyTorch}

  \vspace{0.5em}

  \vspace{10em}

  \cvsection{语言能力}

  \cvtag{中文 (母语)}

  \vspace{0.5em}

  \cvtag{德语 (C1)}

  \vspace{0.5em}

  \cvtag{英语 (B2)}


  %% Yeah I didn't spend too much time making all the
  %% spacing consistent... sorry. Use \smallskip, \medskip,
  %% \bigskip, \vspace etc to make adjustments.
  \medskip

  % use ONLY \newpage if you want to force a page break for
  % ONLY the current column
  \newpage

  %% Switch to the right column. This will now automatically move to the second
  %% page if the content is too long.
  \switchcolumn
  \cvsection{工作经验}

  \cvevent{软件与 AI / 计算机视觉开发工程师}{Innomatik AG}{06/2021 -- 至今}{本斯海姆}

  \vspace{1em}

  % lnux system
  \cvsubsubsection{Linux、开发环境与自动化}

  \vspace{1.0em}

  \begin{itemize}
    \item 使用 Ubuntu WSL 作为主要开发环境，并搭建和配置项目所需的软件及 AI 环境
    \item 使用 SSH 连接公司内部 Linux 服务器并进行开发工作
    \item 编写 Shell 和 Batch 脚本，用于程序启动、参数传递、文件及路径检查，以及所需运行组件的自动配置
  \end{itemize}

  \vspace{1.0em}

  \textbf{技术}：Linux、Ubuntu、WSL、Shell、Bash、Batch、SSH、Git

  \vspace{1.5em}

  % deployment
  \cvsubsubsection{部署、容器化与故障排查}

  \vspace{1.0em}

  \begin{itemize}
    \item 将自主开发的应用程序和 REST 服务部署在本地开发环境、公司内部服务器以及现有企业软件系统中
    \item 使用 Docker 对不同应用程序和 AI 服务进行容器化并运行
    \item 配置 REST 接口、端口及运行参数，使部署后的服务能够在公司内部环境中访问
    \item 独立分析并解决部署和应用运行过程中出现的技术问题，包括 URL 路由错误、端口冲突、Docker 容器无法启动，以及应用程序和配置错误等
  \end{itemize}

  \vspace{1.0em}

  \textbf{技术}：Docker、REST API、FastAPI、ASP.NET Core、Python、C++、C\#

  \vspace{1.5em}

  % software development
  \cvsubsubsection{软件开发与系统集成}

  \vspace{1.0em}

  \begin{itemize}
    \item 使用 Python、C++ 和 C\# 开发模块化软件组件及数据处理 Pipeline，并通过 REST 接口提供服务
    \item 使用 FastAPI 和 ASP.NET Core WebAPI 开发服务，并集成至现有 .NET Blazor Server 系统
    \item 将基于 ONNX 的 AI 组件集成到 C++ 和 C\# 应用程序中，并开发用于内部技术流程的独立桌面应用和控制台程序
  \end{itemize}

  \vspace{1.0em}

  \textbf{技术}：Python、C++、C\#、.NET、ASP.NET Core、Blazor Server、FastAPI、REST API、ONNX

  \vspace{1.5em}

  \newpage

  % database
  \cvsubsubsection{数据库应用与软件开发}

  \vspace{1.0em}

  \begin{itemize}
    \item 在现有企业软件开发过程中持续使用 Microsoft SQL Server
    \item 分析数据库结构、数据类型及数据关系，并理解其在软件模型以及数据读写流程中的映射关系
    \item 创建和修改测试数据，用于验证新增功能及现有软件功能的正确性
  \end{itemize}

  \vspace{1.0em}

  \textbf{技术}：Microsoft SQL Server、SQL、.NET

  \vspace{1.5em}

  % computer vision
  \cvsubsubsection{OCR 与工业图像处理}

  \vspace{1.0em}

  \begin{itemize}
    \item 开发面向真实图像及工程图纸的模块化 OCR Pipeline，包括 Text Detection 与 Text Recognition
    \item 搭建完整的训练与处理流程，包括数据准备、标注、模型训练以及图像预处理和后处理
    \item 将应用部署为本地软件、REST 服务及 Docker 容器化应用
  \end{itemize}

  \vspace{1.0em}

  \textbf{技术}：Python、OpenCV、PyTorch、PaddleOCR、ONNX、FastAPI、Docker

  \vspace{1.5em}


  \cvsection{教育经历}

  \cvevent{土木工程信息技术硕士}{达姆施塔特工业大学}{04/2018 -- 03/2021}{达姆施塔特}

  \vspace{1.0em}

  \cvsubsection{主要方向}
  
  \vspace{1.0em}

  \begin{itemize}
    \item 基于 .NET Framework 的 Revit 插件开发
    \item 使用 SQLite 进行 BIM 数据库开发（CRUD）
    \item 使用 MS Project 进行 BIM 项目管理
  \end{itemize}

  \vspace{1.0em}

  \cvsubsection{科研助理 （HiWi）}

  \vspace{1.0em}

  SCOPE (Semantic Construction Project Engineering)
  \begin{itemize}
    \item 使用 Three.js 开发基于 Web 的 ViewCube
    \item 使用 Protégé 对 IFC 层级结构进行本体建模
  \end{itemize}

  \vspace{1.0em}

  \divider

  \cvevent{桥梁工程学士}{长安大学}{09/2012 -- 06/2016}{西安}

  % \divider

  \medskip


\end{paracol}


\end{document}
